\section{Proof Sketch}
\label{sec:proof-sketch}

The proof begins with a finite one-dimensional threshold game.  An arbitrary
hypothesis is restricted to its predictions on $[N]$, which preserves both
privacy and every risk in the game.  Sampling $n$ rows with replacement from
an empirical database of size $9n$ gives an empirical wrapper whose expected
empirical loss is exactly the finite-game risk.  The corrected
secrecy-of-the-sample calculation of \citet{bunEtAl2015thresholds} preserves
the original privacy parameters uniformly through $\varepsilon=1$.

The homogeneous-set argument of \citet{alonEtAl2018private} is then
rederived at expected loss $1/20$.  A finite coloring records all prediction
marginals, the finite Ramsey theorem of \citet{erdosRado1952} gives a large
homogeneous set, and endpoint privacy comparisons yield a positive adjacent
rise.  Moving one record across the remaining ordered block produces a family
of pairwise private output laws with an off-diagonal margin.  A finite product
experiment and binary-search tree turn this margin into an upper bound on the
homogeneous set.  Comparing the Ramsey lower bound and binary upper bound and
inverting the resulting iterated logarithms gives universal constants
$a_{\rm th},a_\delta>0$ and $N_{\rm th}$ for the one-arm obstruction.

The finite set of private kernels is a compact convex polytope and each
experiment risk is affine.  Finite minimax therefore converts the
algorithm-wise obstruction into a single learner-independent hard prior over
finite threshold experiments.  From this prior we build a hidden-arm kernel:
one uniformly selected arm receives distinct records from an external
$n$-sample, while all other arms are simulated independently.  Fixing all
internal randomness maps an external one-record replacement to zero or one
replacement in the input to $A$, so the kernel inherits the pair
$(\varepsilon_0,\delta_m)$ exactly.

A grand collection of iid arm pools couples the constructed input to an ideal
iid sample from the uniform arm mixture.  The mixture is realized by the
minor-table concept, and Bernstein's inequality bounds the sole overflow
event by $e^{-27/2}$.  The iid latent experiment vector is independent of the
hidden designation, so uniform-arm averaging identifies hidden-arm risk with
mixture risk exactly.  The PAC guarantee and bounded-loss conversion then give
strict prior-average risk below $1/20$, contradicting the hard prior.  The
last arithmetic step transfers the unchanged $\delta_m$ denominator, chooses
$c_\delta=a_\delta/[225(1+\log15)]$ and $c=a_{\rm th}/4$, and retains the
zero-overflow $k=1$ specialization.  Complete derivations appear in the
appendix.
